Este trabajo investiga, desde la práctica, cómo mover esas viejas formas de hacer software hacia métodos más ágiles y modernos, poniendo a prueba si esto ayuda a construir productos que funcionen bien y no se rompan con el tiempo. El problema con los sistemas antiguos es que están tan enredados que cuesta mucho mantenerlos y, lo peor, frenan la innovación, aumentando el riesgo de que todo falle.

Es muy tentador pensar que la solución mágica es borrar todo y arrancar de cero, pero eso casi nunca sirve si antes no cambiamos el "chip" de cómo trabaja el equipo. Por eso, en este artículo no solo proponemos, sino que aplicamos un conjunto de estrategias que realmente funcionan, tomando como nuestro caso central el desarrollo del Portal Agro-comercial del Huila.

La clave para que el software crezca es usar buenas bases (como las reglas SOLID), ser flexibles en el día a día con Scrum, y diseñar un sistema en capas que se pueda partir en pedacitos más manejables. Esto evita que nazca la famosa "deuda técnica". ¿Cómo se logra? Limpiando el código todo el tiempo y dejando que las herramientas automáticas hagan el trabajo pesado, asegurando que el programa siga haciendo lo suyo, pero mejor.

Para cerrar, narramos nuestra experiencia tal cual pasó: desde que entendimos el lío de la metodología tradicional hasta que pusimos en marcha soluciones reales en el Portal Agro-comercial del Huila, como el diseño modular y las pruebas automáticas. Al final, mostramos los resultados y las lecciones que nos quedaron, dejando una guía lista para que otros equipos de ingeniería refactoricen sus procesos y código con éxito.